\chapter{Estado del Arte}
\label{cap:estadoDelArte}

Una vez determinado el objetivo general de este proyecto, continuaremos en este capítulo analizando  
el trasfondo y estado del arte en dos áreas clave: aplicaciones específicas de escalada y mecánicas de gamificación 
aplicadas en otros deportes.

Se presenta primero un análisis del panorama tecnológico actual examinando aplicaciones relevantes, sus funcionalidades y debilidades. 
Posteriormente se definen los principios de diseño y mecánicas de gamificación que se pueden extrapolar para el desarrollo de la aplicación.

\section{Análisis de aplicaciones de escalada}

En esta sección se analizan las principales aplicaciones que se emplean para registrar la actividad en la escalada. Las aplicaciones 
se han agrupado en base a su proposito definido.

\subsection{Aplicaciones integradas con el rocódromo}

Estas aplicaciones tienen una integración estrecha con los rocódromos, actuando como una herramienta para el registro de rutas y sesiones.

\paragraph{TopLogger}\citep{TopLogger}:
TopLogger es una aplicación móvil que permite registrar el estado con las rutas de los rocódromos asociados y facilita la identificación de estas por 
ubicación, color y grado. Para determinar las posiciones de las rutas en un rocódromo realiza a través de un mapa cenital que permite moverte 
y seleccionar las rutas deseadas rápidamente. En el perfil se pueden visualizar numerosas estadísticas sobre datos de las sesiones de escalada y para 
su capa de gamificación tienen implementado un sistema de niveles y clasificaciones por rocódromo. 

Entre sus  limitaciones la primera que salta a la vista es la sobrecarga de información en cada pantalla, haciendo poco intuitiva y 
poco accesible la aplicación para pefiles con un nivel tecnólogico bajo. En cuanto a la gamificación, al hacer las clasificaciones por rocódromo se puede 
engañar a la aplicación y posicionarte primero, consiguiendo el efecto adverso de motivación para los escaladores legítimos. Además de esto, cuenta con poca 
integración con los rocódromos en España y con una capa de pago si tu rocodromo cuenta con muchos escaladores habituales.

\paragraph{WeClimb}\citep{WeClimb}:
WeClimb presenta funcionalidades similares a TopLogger, incluyendo un registro de rutas por color y grado pero difiriendo en la posición de las rutas.
En esta aplicación para seleccionar la ruta usa un sistema de tarjetas de manera que seleccionando una sección de un mapa, se muestra 
una serie de tarjetas con las fotos de cada ruta facilitando el reconocimiento visual. En las imagenes de las rutas existe un mecanismo 
para poder resaltar las presas(agarres) que se pueden usar, dando la posibilidad de crear retos personalizados. Además de esto la aplicación 
cuenta con una pestaña para eventos, un historial con las tarjetas de las rutas completadas y clasificaciones por rocódromo. 

Como limitaciones encontramos que su despliegue está centrado para Francia por lo que no hay centros ni traducción de la app. Actualmente carece 
de una capa de gamificación profunda, a excepción de las clasificaciones, pero que como en TopLogger se pueden engañar creando un efecto adverso. 
Por último el sistema de tarjetas para las rutas puede ser algo tedioso si hay numerosas rutas en una zona haciendo lento el proceso de selección 
de estas.


\subsection{Aplicaciones centradas en la comunidad}

En este grupo se prioriza la conexión entre escaladores frente al registro de todos los ascensos. Funcionan como redes sociales especializadas 
para compartir los ascensos más importantes.

\paragraph{KAYA\citep{KAYA}}
KAYA se posiciona como una red social para escaladores (especialmente orientada a escalada en roca natural), centrada en compartir 
ascensos con formato de publicaciones, interactuar con usuarios y seguir la progresión de otros. Las publicaciones pueden ser registros de 
una ruta, registro de una sesión de escalada o subir un video corto en formato vertical de una solución. En cuanto a elementos gamificación 
tiene un sitema de Badges en base al número de rutas registradas en la aplicación motivando a los escaladores a registrar más rutas. 

Como limitaciones determinamos que la aplicación está más orientada al registro de rutas relevantes para el escalador más que el registro completo 
de toda su actividad. Además, la aplicación esta principalmente centrada en escalada en roca natural por lo que tiene poca integaración con rocódromos.


\subsection{Aplicaciones con hardware específico}

Existen soluciones que combinan software con elementos físicos del rocódromo para crear experiencias de entrenamiento más interactivas.

\paragraph{MoonBoard\citep{MoonBoard}}

MoonBoard es una instalación física de entrenamiento que combina un muro de escalada con una aplicación. La particularidad de este sistema reside
en que la posición de las presas y la inclinación del muro son identicas en todas sus instalaciones, permitiendo así una experiencia común en cualquier
centro de escalada. La aplicación se conecta al muro y permite elegir entre miles de rutas que los escaladores han creado y una vez seleccionado uno,
las presas que se pueden usar se iluminan. MoonBoard tiene un sistema de ranking global que funciona por un sistema de ascensos validados por la comunidad permitiendo a los usuarios competir 
con escaladores de otros paises y visualizar el progreso en comparación con la media mundial.

Como limitaciones destacamos que MoonBoard es un sistema diseñado para entrenamientos de alta instesidad por lo que no es accesible a escaladores novatos. 
Además de esto MoonBoard es un sistema cerrado por lo que únicamente se pueden regsitrar las rutas completadas en el funcionando como una herrmamienta
de registro parcial. 

\section{Análisis de aplicaciones de otros deportes}

En esta sección se examinan mecánicas de gamificación probadas en otras disciplinas y aplicaciones de fitness con el objetivo de 
identificar ideas transferibles al contexto del rocódromo.

\subsection{Competiciones y ranking por segmentos (Strava)}

\paragraph{Strava\citep{Strava}}

Strava es una aplicación de registro de sesiones de corredores y ciclistas. Parte de su exito viene de su concepto de segmentos, que son tramos
específicos de una ruta donde los usuarios registran y comparan los tiempos que han tardado en completarlo mediante GPS. Este sistema transforma
un entrenamiento indicvidual en una competición virtual y el modelo se puede extrapolar a la escalada tratando cada vía o bloque como un segmento,
donde se pueden medit el número de intentos o tiempo. En cuanto a la gamificación tiene un sistema de Badges que son títulos como KOM/QOM (King/Queen of 
the Mountain) que fomentan la motivación de los usuarios.

Como limitaciones en el enfoque de la escalada determinamos que un sistema de segmentos basados en tiempo en la escalada es contraproducente ya que
se puede comprometer la tecnica y seguridad del ascenso.

\subsection{Coleccionismo de logros (Pokémon Go)}

\paragraph{Pokémon Go\citep{PokemonGO}}
Pokémon Go es una aplicación del universo Pokémon centrada en la gamificación del desplazamiento físico con la vinculación de recompensas virtuales
y el progreso en la narrativa de un juego. Este sistema permite convertir el esfuerzo en un incentivo para obtener el sentimiento de coleccionsimo 
y descubrimiento. La aplicación en un ejemplo de la implementación del sistema PBL (Points Badges and Leaderboards) previamente definido. 

La desventaja principal de esta aplicación es que el objetivo del usuario deja de ser el deporte y pasa a ser la recompensa del juego, 
creando una posible dependencia negativa en los usuarios.


\section{Conclusiones del Estado del Arte e identificación de la oportunidad}
 
Para finalizar este capítulo, se presentan las concolusiones obtenidas tras el análisis comparativo entre las aplicaciones, 
permitiendo identificar la necesidad del mercado y la propuesta de valor de nuestra aplicación.

\subsection{Síntesis de la problemática actual}

Tras analizar las soluciones existentes, se observa que el mercado de aplicaciones de escalada presenta problemas en tres aspectos críticos:

\begin{itemize}
\item \textbf{Experiencia de Usuario (UX) deficiente:} Aplicaciones como TopLogger o WeClimb ofrecen una gran densidad de datos, 
pero a costa de interfaces saturadas que dificultan el uso y registro de rutas en la aplicación. La barrera tecnológica aleja 
al escalador casual que solo busca un registro sencillo de su progresión.

\item \textbf{Gamificación superficial:} La mayoría de las herramientas se limitan a clasificaciones numéricas (leaderboards) 
fácilmente manipulables, lo que genera desmotivación. Falta un sistema de recompensas basado en el esfuerzo, la constancia y el "coleccionismo" 
de hitos personales que no dependa únicamente de ser el mejor del ranking.

\item \textbf{Dependencia de hardware o geografía:} Existen buenas soluciones como MoonBoard, pero están limitadas a sus muros de escalada específicos, 
o aplicaciones como WeClimb y TopLogger cuya implantación en España es inexistente. No existe una herramienta universal, intuitiva y "gamificada" 
que el escalador pueda llevar de un rocódromo a otro.
\end{itemize}

\subsection{Identificación de la oportunidad}

El análisis de aplicaciones como \textbf{Strava} y \textbf{Pokémon Go} revela que existe un gran potencial en la adaptación de mecánicas de gamificación
de otros sectores a la escalada. La oportunidad reside en desarrollar una solución que combine:

\begin{enumerate}
\item \textbf{Registro visual y ágil:} Adoptar un sistema híbrido de mapas centiales dinámicos y sistema de tarjetas para minimizar el tiempo 
de interacción de escalador para reconocer y acceder a la ruta con el móvil durante la sesión.
\item \textbf{Gamificación basada en PBL:} Realizar una implementación correcta de el sistema PBL basado en el modelo RAMP de forma que las estadísticas
funcionen como los \textit{Points} para mantener la motivación y constancia de los usuarios, además de la  implementación de Leaderboards entre 
grupos de amigos eliminando el factor de los abusos de los sistemas de posicionamiento.

\item \textbf{Enfoque social y comunitario:} Fomentar la interacción mediante personalización del perifl y la comparación de estadísticas 
personales, creando un sentimiento de pertenencia a la comunidad del rocódromo local.
\end{enumerate}

En conclusión, este proyecto se posiciona entre als herramientas de entrenamiento puro y las redes sociales generales, 
ofreciendo un sistema de entrenamiento que registra datos de las sesiones y motiva activamente al  escalador a través de una 
experiencia lúdica y bien diseñada.